%\documentclass[a4paper, 10pt, twocolumn, journal, twoside]{IEEEtran}
\documentclass[a4paper, 12pt, onecolumn, journal, twoside, draftclsnofoot]{IEEEtran}

\newcommand*{\DOUBLECOL}{}	% turns this on for double column
\newlength{\subfigureheight}
\newlength{\figurewidth}
\ifdefined\DOUBLECOL
\setlength{\subfigureheight}{2.45in}
\setlength{\figurewidth}{3.5in}
\else
\setlength{\subfigureheight}{2.35in}
\setlength{\figurewidth}{6.5in}
\fi
\usepackage{amsfonts,amsmath,amssymb,bm,cancel,color,comment,epsfig,enumerate,empheq,floatrow,hyperref,mathdots,mathrsfs}
\usepackage{placeins}%\usepackage{algorithmic}
\usepackage[vlined,ruled,linesnumbered]{algorithm2e}
\usepackage{graphicx}
%\usepackage{amsthm}
\usepackage[super]{nth}
\usepackage[noadjust]{cite}
\usepackage{subcaption}
\usepackage{tabularx}
\usepackage{tabularray}		% make thing horizontal lines in tabular
\UseTblrLibrary{booktabs}	% make thing horizontal lines in tabular

\graphicspath{{eps/}}

\hypersetup{
	colorlinks=true,
	linkcolor=red,
	citecolor=green,
	filecolor=magenta,
	urlcolor=blue
}

\newenvironment{Ualgorithm}[1][htpb]{\def\@algocf@post@ruled{\kern\interspacealgoruled\hrule  height\algoheightrule\kern3pt\relax}%
\def\@algocf@capt@ruled{under}
\begin{algorithm}[#1]}
{\end{algorithm}}

\setcounter{MaxMatrixCols}{30}
\DeclareSymbolFont{grb}{OML}{cmm}{b}{it}
\DeclareMathSymbol{\zetab}{\mathord}{grb}{"10}
\DeclareMathSymbol{\etab}{\mathord}{grb}{"11}
\DeclareMathSymbol{\thetab}{\mathord}{grb}{"12}
\DeclareMathSymbol{\kappab}{\mathord}{grb}{"14}
\DeclareMathSymbol{\lambdab}{\mathord}{grb}{"15}
\DeclareMathSymbol{\mub}{\mathord}{grb}{"16}
\DeclareMathSymbol{\nub}{\mathord}{grb}{"17}
\DeclareMathSymbol{\rhob}{\mathord}{grb}{"1A}
\DeclareMathSymbol{\sigmab}{\mathord}{grb}{"1B}
\DeclareMathSymbol{\taub}{\mathord}{grb}{"1C}
\DeclareMathSymbol{\phib}{\mathord}{grb}{"1E}
\DeclareMathSymbol{\psib}{\mathord}{grb}{"20}
\DeclareMathSymbol{\omegab}{\mathord}{grb}{"21}
\DeclareMathSymbol{\epsilonb}{\mathord}{grb}{"22}
\DeclareMathSymbol{\varphib}{\mathord}{grb}{"27}
\input{Symbol_Shortcut.tex}

\SetKwRepeat{Do}{do}{while}%

\newtheorem{theorem}{\bf Theorem}		% by Victor
\newtheorem{corollary}{\bf Corollary}	% by Victor
\newtheorem{definition}{\bf Definition}	% by Victor

%\centerfigcaptionstrue

%\pagestyle{empty}

\makeatletter
\def\ifundefined{\@ifundefined}
\makeatother

\setlength{\parskip}{0ex}

% generating the double accent on o, i.e. \fH{o}, e.g. Erdos Renyi 
\newcommand\fH[1]{\sbox0{#1}\dimen0=\ht0 \advance\dimen0 -1ex
  \sbox2{\'{}}\sbox2{\raise\dimen0\box2}%
  {\ooalign{\hidewidth\kern.1em\copy2\kern-.5\wd2\box2\hidewidth\cr\box0\crcr}}}

\title{\vspace{-0.2em} \huge Model-Driven Neural Network Based MIMO Channel Estimator}

\ifdefined \BLIND
\author{Carrson C. Fung $^1$ \\ %and Author 2$^2$ \\
$^1$ NYCU  \\
%email: \emph{Author 1's email, Author 2's email} \\
\thanks{Funding sources.}
}
\else
\author{Shanik Hubert, ~\IEEEmembership{Student Member,~IEEE}, and Carrson C. Fung$^*$,~\IEEEmembership{Member,~IEEE} \\
%Institute of Electronics, College of ECE, National Chiao Tung University, Hsinchu, Taiwan 300 %\\
%Voice: +886-3-573-1862$^2$ Fax: +886-3-572-4361$^2$ email: \emph{rusly.eic04g@nctu.edu.tw, c.fung@ieee.org \vspace{-1.5em}} 
\thanks{This work has been partially supported by Google exploreCSR Grant 111Q90004C.

Shanik Hubert and Carrson C. Fung  are both affiliated with the Department of Electrical and Computer Engineering, National Yang Ming Chiao Tung University, Hsinchu 300, Taiwan. (e-mail: dmytro.ivakhnenkov@gmail.com, c.fung@ieee.org)}
}
\fi


\markboth{IEEE Trans. on Communications}{Latency Minimization for Uplink RSMA System}
\begin{document}
\maketitle

%******************************************************************
\begin{abstract}
%******************************************************************
\end{abstract}

\begin{IEEEkeywords}
\end{IEEEkeywords}

%%%%%%%%%%%%%%%%%%%%%%%%%%%%%%%%%%%%%%%%%%%%%%%%%%%%%%%%%%%%%%%%%%
\section{Introduction}	\label{intro}
%%%%%%%%%%%%%%%%%%%%%%%%%%%%%%%%%%%%%%%%%%%%%%%%%%%%%%%%%%%%%%%%%%


  
%%%%%%%%%%%%%%%%%%%%%%%%%%%%%%%%%%%%%%%%%%%%%%%%%%%%%%%%%%%%%%%%%%
\section{Methodology} \label{sec_meth}
%%%%%%%%%%%%%%%%%%%%%%%%%%%%%%%%%%%%%%%%%%%%%%%%%%%%%%%%%%%%%%%%%%

\subsection{Uplink Latency Minimization}
%******************************************************************
For MIMO uplink (UL), the receive signal is written as
\begin{equation*}
	\y = \sum_{k=1}^K  \H_k \F_k \s_k + \w \in \complex^{N_R},
\end{equation*}
where 
$\H_k \in \complex^{N_R \times N_k}$ denotes the uplink channel between the $k$th user and the AP, $\F_k \in \complex^{N_k \times d_k}$ denotes the precoding matrix for the $k$th user, $d_k$ denotes the number of spatial streams user $k$ is sending to the AP (and $d_k \leq N_k$), $\s_k \in \complex^{d_k}$ denotes the transmitted signal for the $k$th user, $\w \in \complex^{N_R}$  denotes the AWGN, with $N_k$ denoting the number of antennas at the $k$th user and $N_R$ denotes the number of antennas at the AP.  It is assumed that $E \left[ \s_k \s_k^H \right] = \I_{N_k}$.  Let 
\begin{equation*}
	x_k = \F_k \s_k,
\end{equation*}
so
\begin{align*}
	E \left[ \left| \x_k \right\|_2^2 \right] &= \F_k E \left[ \s_k \s_k^H \right] \F_k^H \\
	&= \F_k \F_k^H  \\
	&=  \left\| \F_k \right\|_F^2  \leq P_k,
\end{align*}
where $P_k$ is the transmit power from the $k$th user.  If SIC is used at the receiver, then the achievable rate at the receiver is
\begin{equation*}
	R_{sum} = \sum_{k=1}^K  \log_2 \det \left(  \I_{N_R}  +  \H_k \F_k     \left(  \sigma_w^2 \I_{N_R}  + \sum_{j \neq k}  \H_j \F_j \F_j^H \H_j^H  \right)^{-1}     \F_k^H \H_k^H  \right) 
\end{equation*}

Then it seems we can formulate the latency minimization problem as
\begin{equation}
	\begin{split}
		\min_{\F_k, B_k, \forall k}  &\;  \sum_{k=1}^K  \frac{ B_k } { R_{sum}  }  \\
		s.t. &\;  \left\| \F_k \right\|_F^2  \leq P_k,  \forall k, \\
		&\; B_k \in \left\{  B_{k,\min}, \ldots, B_{k,\max} \right\}.
	\end{split}
\label{eq_prob_form_basic}
\end{equation}
Possible extension include optimizing not with the sum rate, but using rate such as
\begin{enumerate}[1.]
	\item Max-min rate: 
	\begin{equation*}
		\max_{\F_k, \forall k}  \min_k  R_k.
	\end{equation*}
This has a strong fairness guarantee, often leads to equal rates across all users if we can find a feasible solution.  This is the worst-case that Carrson was talking about.
	\item The proportional fairness:
	\begin{equation*}
		\max_{\F_k, \forall k}   \sum_{k=1}^K  \log (\R_k), 
	\end{equation*}
which balances the throughput and fairness. 
	\item Weighted fairness 
	\begin{equation*}
		\max_{\F_k, w_k, \forall k}  w_k \log (R_k), 
	\end{equation*}
where $w_k$ are weights for QoS differentiation, i.e., higher $w_k$ is, more resource will be allocated the $k$th user.
\end{enumerate}


\subsection{Uplink Latency Minimization Using RSMA}
%******************************************************************
 
\begin{comment}
%***************************************************************************
\begin{figure}[ht]
    \centering
	\subcaptionbox{Indoor factory (InF-DL) channel at 2.5 GHz. \label{fig_NMSE_vs_snr_InF_2_5_ISTA_Net}} {
		\includegraphics[width = 0.47\columnwidth]{./figs/ISTA_LS_Net_eigenmode_CNN03_vs_ISTA_Net_InF_DL_2_5.png}
	} %
	\subcaptionbox{Urban Macro (UMa) channel at 28 GHz.    \label{fig_NMSE_vs_snr_UMa_28_ISTA_Net}} { 
		\includegraphics[width = 0.47\columnwidth]{./figs/ISTA_LS_Net_eigenmode_CNN03_vs_ISTA_Net_UMa_28.png}
	}
    \caption{Normalized mean-squared error (NMSE) vs. SNR (both in dB) results for MIMO channel estimation for our proposed ISTA-LS-Net eigenmode with CNN with L = 5 layers, least squares (LS), linear minimum mean-squared estimator (LMMSE) and ISTA-Net \cite{liu2021} with $L = 10$ layers.  $n_T = 36$, $n_R = 4$, $T = 36$.  The results are plotted using the InF-DL 2.5 GHz and UMa channels at 28 GHz as stated in  \cite{PL_for}.}
    \label{fig_ISTA_LS_Net}
\end{figure}   
%***************************************************************************
\end{comment}






%%%%%%%%%%%%%%%%%%%%%%%%%%%%%%%%%%%%%%%%%%%%%%%%%%%%%%%%%%%%%%%%%%
\section{Simulation Results} \label{sec_simu}
%%%%%%%%%%%%%%%%%%%%%%%%%%%%%%%%%%%%%%%%%%%%%%%%%%%%%%%%%%%%%%%%%%




\begin{comment}
%*************************************************************************
\begin{table}[!htb]
	\small
	\caption{Simulation parameters for PN-IEKF.}
	\centering
	\begin{tabularx}{\textwidth}{Xllll}
		\hline
  		\multicolumn{4}{c}{Simulation Parameters}\\%{Precoder Design Parameters} \\
		\hline
		Number of vertices $(M)$ & 10 & 10 & 30 \\
		SNR & 5 dB & 20 dB & 20 dB\\
		Variance of process noise $\widebar{\u}[n]$ ($\sigma_u^2$) & 1 & 1 & 1 \\
		Variance of process noise $\u_\omega[n]$ ($\sigma_\omega^2$) & 0.045 & 0.002 & 0.001\\
		Sampling period $(T)$ & 1 & 1 & 1\\
		Number of Monte Carlo experiments & 100 & 100 & 100 \\
		Number of time samples for Warm start & 90 & 90 & 90\\
		Max. number of IEKF iterations (MaxIter) & 3 & 4 & 6 \\
		Max. number of ISTA iterations & 15000 & 15000 & 15000\\
		Transition time $(n_0)$ & 14000 & 14000 & 14000\\
		Number of time samples used in $T$-squared test $(N)$& 200 & 200 & 200\\
		Sparsity regularization factor $(\beta)$ & 0.18 & 0.029 & 0.01\\
		Edge probability parameter $(p)$ & 0.4 & 0.4 & 0.4\\
		False alarm rate $(\alpha)$ & $10^{-4}$ & $10^{-5}$ & $10^{-5}$\\
		Cool down time samples & 6000 & 6000 & 6000 \\
		Initial MSE of estimated graph signals $(a)$ & 100 & 100 & 100 \\
		Initial MSE of estimated graph $(b)$ & 1 & 1 & 1 \\
		\hline
	\end{tabularx}
	\label{tab_sim_params_pn_iekf}
\end{table}
%*************************************************************************
%**************************************************************************
\begin{figure}[htp]
    \centering
    \subcaptionbox{NMSD vs. time sample $n$. \label{fig_NMSD_snr20dB_M10}} {
		%\includegraphics[width=0.47\columnwidth]{eps/PN_IEKF/NMSD_RE/20dB_10node_all_NMSD.eps}
	}%
	\subcaptionbox{RE vs. time sample $n$. \label{fig_RE_snr20dB_M10}} {
		%\includegraphics[width=0.47\columnwidth]{eps/PN_IEKF/NMSD_RE/20dB_10node_all_RE.eps}
	}
\caption{NMSD and RE vs. time sample $n$ comparison between the oracle, the PN-IEKF method, three LMS approaches proposed in \cite{vlaski2018}. $M=10$, $\text{SNR} = 20 \text{ dB}$. Transition occurs at $n_0 = 14000$ sample.}
\label{fig_NMSD_RE_snr20dB_M10}
\end{figure}
%**************************************************************************


\subsection{Simulation Results of the PN-IEKF-VAR Using the VAR Model}
%******************************************************************

\end{comment}

%%%%%%%%%%%%%%%%%%%%%%%%%%%%%%%%%%%%%%%%%%%%%%%%%%%%%%%%%%%%%%%%%%
\section{Conclusion}		\label{sec_concl}
%%%%%%%%%%%%%%%%%%%%%%%%%%%%%%%%%%%%%%%%%%%%%%%%%%%%%%%%%%%%%%%%%%


%****************************************
\begin{thebibliography}{25}


\bibitem{chang2015semm}
C.-Y. Chang and C.C. Fung, ``Sparsity enhanced mismatch model for robust spatial intercell interference cancelation in heterogeneous networks,'' \emph{IEEE Trans. on Communications}, vol. 63(1), pp. 125-139, Jan. 2015.

\bibitem{daubechies2004}
I. Daubechies, M. Defrise and C. De Mol, ``An iterative thresholding algorithm for linear inverse problems with a sparsity constraint,'' \emph{Communications on Pure and Applied Mathematics}, vol. 57(11), pp. 1413-1457, 2004.

\bibitem{beck2009}
A. Beck and M. Teboullen, ``A fast iterative shrinkage-thresholding algorithm for linear inverse problems,'' \emph{SIAM Journal on Imaging Sciences}, vol. 2(1), pp. 183-202, 2009.

\bibitem{parikh2013}
N. Parikh and S. Boyd, ``Proximal algorithms,'' \emph{Foundations and Trends in  Optimization}, vol. 1(3), pp. 123-231, 2013.

\bibitem{cohen2019}
R. Cohen \emph{et al.}, ``Deep convolutional robust PCA with application to ultrasound imaging,'' \emph{Proc. of the IEEE Intl. Conf. on Acoustics, Speech and Signal Processing}, pp. 3212-3216, Brighton, UK, May 2019.

\bibitem{soloman2019}
O. Soloman \emph{et al.}, ``Deep unfolded robust PCA with application to clutter suppression in ultrasound,'' \emph{IEEE Trans. on Medical Imaging}, vol. 39(4), pp. 1051-1063, Sep. 2019.

\bibitem{heath2016}
R.W. Heath Jr., N. Gonz\'alez-Prelcic, S. Rangan, W. Roh and A.M. Sayeed, ``An overview of signal processing techniques for millimeter wave MIMO systems,'' \emph{IEEE Journal of Selected Topics in Signal Processing}, vol. 10(3), pp. 436-453, Dec. 2016.



\bibitem{kay1993}
S. Kay, \emph{Statistical Signal Processing: Estimation Theory}, Prentice Hall, 1993.

\bibitem{eisen2019}
M. Eisen, C. Zhang, L.F.O. Chamon, D.D. Lee and A. Ribeiro, ``Learning optimal power allocations in wireless systems,'' \emph{IEEE Trans. on Signal Processing}, vol. 67(10), pp. 2775-2790, May 2019.

\bibitem{nesterov2017}
Y. Nesterov and V. Spokoiny, ``Random gradient-free minimization of convex functions,'' \emph{Foundations of Computational Mathematics}, vol. 17(2), pp. 527-566, 2017.

\bibitem{sutton2000}
R.S. Sutton, D.A. McAllester, S.P. Singh and Y. Mansour, ``Policy gradient methods for reinforcement learning with function approximation,'' \emph{Proc. of the Advances in Neural Information Processing Systems}, pp. 1057-1063, 2000.

\bibitem{eisen2020}
M. Eisen and A. Ribeiro, ``Optimal wireless resource allocation with random edge graph neural networks,'' \emph{IEEE Trans. on Signal Processing}, vol. 68, pp. 2977-2991, 2020.

\bibitem{PL_for}
3GPP, ``Section 7: Channel model(s) for 0.5-100 GHz'' \emph{5G Study on channel model for frequencies from 0.5 to 100 GHz (3GPP TR 38.901 version 16.1.0 Release 16)}, 2020.

\bibitem{liu2021}
S. Liu and X. Huang, ``Sparsity-aware channel estimation for mmWave massive MIMO: A deep CNN-based approach,'' \emph{China Communications}, Jun. 2021.






\begin{comment}
\bibitem{bridgeford2022}
E.W. Bridgeford, A. Loftus and J.T. Vogelstein, \emph{Hands-on Network Machine Learning with Scikit-Learn and Graspolgic}. [Online]. Available: \href{http://docs.neurodata.io/graph-stats-book/appendix/ch12/rdpgs.html}{\textit{http://docs.neurodata.io/graph-stats-book/appendix/ch12/rdpgs.html}}.

% cvx
\bibitem{grant}
M. Grant and S. Boyd, \emph{CVX: Matlab software for disciplined convex programming}, \href{http://cvxr.com/cvx/}{\textit{http://cvxr.com/cvx/}}.

% YALMIP
\bibitem{yalmip}
%YALMIP, ``Yet Another Linear Matrix Inequality Parser,'', \textcolor{blue}{\url{http://users.isy.liu.se/johanl/yalmip/}}.
YALMIP, ``Yet Another Linear Matrix Inequality Parser,'', \href{http://users.isy.liu.se/johanl/yalmip/}{\textit{http://users.isy.liu.se/johanl/yalmip/}}.

% Graph toolbox
\bibitem{girault2017}
B. Girault, S. S. Narayanan, A. Ortega, P. Gonalves, and E. Fleury, ``Grasp: A matlab toolbox for graph signal processing,'' \emph{IEEE International Conference on Acoustics, Speech and Signal Processing}, pp. 6574-6575, Mar. 2017.  Available: \emph{\textcolor{blue}{https://gforge.inria.fr/projects/grasp}}.

\bibitem{defferrard}
M. Defferrard, L. Martin, R. Pena, and N. Perraudin, ``Pygsp: Graph signal processing in python''. Online. \href{https://github.com/epfl-lts2/pygsp}{\textit{https://github.com/epfl-lts2/pygsp}}.


% Graph neural net
\bibitem{bruna2014}
J. Bruna, W. Zaremba, A. Szlam, and Y. LeCun, ``Spectral networks and locally connected networks on graphs,'' \emph{Proc. of the Intl. Conf. on Learning Representations}, Banff, AB, Canada, Apr. 2014.

\bibitem{kipf2017}
T.N. Kiph and M. Welling, ``Semi-supervised classification with graph convolutional networks,'' \emph{Proc. of the International Conf. on Learning Representations}, 2017.

\bibitem{du2018TAGCN}
J. Du \emph{et al.}, ``Topology adaptive graph convolutional networks,'' \emph{arXiv:1710.10370v5}, Feb. 2018.


\bibitem{bronstein2017}
M.M. Bronstein, J. Bruna, Y. LeCun, A. Szlam, and P. Vandergheynst, ``Geometric deep learning,'' \emph{IEEE Signal Processing Magazine}, vol. 34(4), pp. 18-42, Jul. 2017.

%\bibitem{shi2019}
%J. Shi and J.M.F. Moura, ``Graph signal processing: Modulation, convolution, and %sampling,'' \emph{arXiv:1912.06762}, Dec. 2019.


\bibitem{yuan2006}
M. Yuan and Y. Lin, ``Model selection and estimation in regression with grouped variables,'' \emph{Journal of the Royal Society. Series B,} vol 68(1), pp. 49-67, 2006.

\bibitem{bassett2011}
D.S. Bassett, N.F. Wymbs, M.A. Porter, P.J. Mucha, J.M. Carlson, and S.T. Grafton, ``Dynamic reconfiguration of human brain networks during learning,'' \emph{Proceedings of the National Academy of Sciences}, vol. 108(18), pp. 7641-7646, May 2011.

\bibitem{cressie2011}
N. Cressie and C.K. Wilke, \emph{Statistics for Spatio-Temporal Data}, Hoboken, NJ, USA, Wiley, 2011.

\bibitem{condat2013}
L. Condat, ``A primal-dual splitting method for convex optimization involving Lipschitzian, proximable and linear composite terms,'' \emph{Journal of Optimization Theory and Applications}, vol. 158(2), pp. 460-479, 2013.

\bibitem{combettes2012}
P.L. Combettes and J.C. Pesquet, ``Primal-dual splitting algorithm for solving inclusions with mixtures of composite, Lipschitzian, and parallel-sum type monotone operators,'' \emph{Set-Valued and Variational Analysis}, vol. 20(2), pp. 307-330, 2012.

\bibitem{komodakis2015}
N. Komodakis and J.-C. Pesquet, ``Playing with duality: An overview of recent primal-dual approaches for solving large-scale optimization problems,'' \emph{IEEE Signal Processing Magazine}, vol. 32(6) pp. 31-54, Nov. 2015.

\bibitem{liu2019}
Y. Liu \emph{et al.}, ``Graph learning based on spatiotemporal smoothness for time-varying graph signal,'' \emph{IEEE Access}, vol. 7, pp. 62372-62386, May 2019.

\bibitem{wu2023}
Z.-Y. Wu and C.C. Fung, ``Online graph topology estimation using proximal Newton iterated Kalman filtering,'' \emph{Under preparation}.

\bibitem{lee2014}
J.D. Lee, Y. Sun and M.A. Saunders, ``Proximal Newton-type methods for minimizing composite functions,'' \emph{SIAM Journal on Optimization}, vol. 24(3), pp. 1420-1443, 2014.  

\bibitem{jazwinski1970}
A.J. Jazwinski, \emph{Stochastic processes and filtering theory}, Academic Press, New York, 1970.

\bibitem{bell1993}
B.M. Bell and F.W. Cathey, ``The iterated Kalman filter update as a Gauss-Newton method,'' \emph{IEEE Trans. on Automatic Control}, vol. 38(2), pp. 294-297, Feb. 1993.

\bibitem{bell1994}
B.M. Bell, ``The iterated Kalman smoother as a Gauss-Newton method,'' \emph{SIAM Journal on Optimization}, vol. 4(3), pp. 626-636, Aug. 1994.

\bibitem{skoglund2015}
M.A. Skoglund, G. Hendeby and D. Axehill, ``Extended kalman filter modifications based on an optimization view point,'' \emph{18th Intl. Conf. on Information Fusion}, pp. 1856-1861, 2015.

\bibitem{parikh2013}
N. Parikh and S. Boyd, ``Proximal algorithms,'' \emph{Foundations and Trends in Optimization}, vol. 1(3), pp. 123-231, 2013.

\bibitem{evap}
\emph{California Irrigation Management Information System}. [Online].  Available: \href{https://cimis.water.ca.gov/}{\textit{https://cimis.water.ca.gov/}}

\bibitem{adni}
\emph{Alzheimer's Disease Neuroimaging Inititative}. [Online]. Available:\href{http://adni.loni.usc.edu/}{\textit{http://adni.loni.usc.edu/}}

\bibitem{openneuro}
\emph{OpenNEURO}.  [Online]. Available: \href{https://openneuro.org/}{\textit{https://openneuro.org/}}

\bibitem{bielczyk2019}
N. Z. Bielczyk, S. Uithol, T. van Mourik, P. Anderson, J.C. Glennon, J.K. Buitelaar, ``Disentangling causal webs in the brain using functional magnetic resonance imaging: A review of current approaches,'' \emph{Network Neuroscience}, vol. 3(2), pp. 237-273, 2019.

\bibitem{baingana2013}
B. Baingana, G. Mateo, and G.B. Giannakis, ``Dynamical structural equation models for tracking topologies of social networks,'' \emph{Proc. of Workshop on Computational Advances in Multi-sensor Adaptive Processing}, Saint-Martin, pp. 292-295, Dec. 2013.

\bibitem{baingana2015}
B. Baingana and G.B. Giannakis, ``Switched dynamic structural equation models for tracking social network topologies,'' \emph{Proc. of the IEEE Global Conf. on Signal and Information Processing}, Orlando, FL, USA, Dec. 2015.

\bibitem{baingana2017}
B. Baingana and G.B. Giannakis, ``Tracking switched dynamic network topologies from information cascades,'' \emph{IEEE Trans. on Signal Processing}, vol. 65(4), pp. 985-997, 2017.

\bibitem{shen2017dyn}
Y. Shen, B. Baingana, and G.B. Giannakis, ``Tensor decompositions for identifying directed graph topologies and tracking dynamic networks,'' \emph{IEEE Trans. on Signal Processing}, vol. 65(14), pp. 3675-3687, Jul. 2017.

\bibitem{bishop2006}
C.M. Bishop, \emph{Pattern recognition and Machine Learning}, in Information Science and Statistics, New York, NY, USA, Springer-Verlag, 2006.

\bibitem{harring2012}
J. Harring, B. Weiss, and J.-C. Hsu, ``A comparison of methods for estimating quadratic effects in nonlinear structural equation models,'' \emph{Psychological Methods}, vol. 17, pp. 193-214, 2012.

\bibitem{finch2015}
W.H. Finch, ``Modeling nonlinear structural equation models: A comparison of the two-stage generalized additive models and the finite mixture structural equation model,'' \emph{Structural Equation Modeling: A Multidisciplinary Journal}, vol. 22(1), pp. 60-75, 2015.

\bibitem{lim2015}
N. Lim, F. D'Alche-Buc, C. Auliac, and G. Michailidis, ``Operator-value kernel-based vector autoregressive models for network inference,'' \emph{Machine Learning}, vo. 99(3), pp. 489-513, 2015.


% Graph neural net
\bibitem{scarselli2009}
G. Scarselli, M. Gori, A.C. Tsoi, M. Hagenbuchner, and G. Monfardini, ``The graph neural network model,'' \emph{IEEE Trans. on Neural Networks}, vol. 20(1), pp. 61-80, Jan. 2009.


%\bibitem{shafipour2019}
%R. Shafipour, S. Segarra, A.G. Marques, G. Mateos, ``Identifying the topology of undirected networks from diffused non-stationary graph signals,'' \emph{arXiv:1801:03862v2}, Jan. 2019.


%\bibitem{boyd2010}
%S. Boyd, N. Parikh, E. Chu, B. Peleato, and J. Eckstein, ``Distributed optimization and statistical learning via the alternating direction of method of multipliers,'' \emph{Foundations and Trends in Machine Learning}, vol. 3(1), pp. 1-122, 2010.

     
%\bibitem{girault2015}
%B. Girault, P. Goncalves, and E. Fleury, ``Translation on graphs: an isometric shift operator,'' \emph{IEEE Signal Processing Letters}, vol. 22(12), pp. 2416-2420, Dec. 2015.

     
%\bibitem{bruna2014}
%J. Bruna, W. Zaremba, A. Szlam, and Y. LeCun, ``Spectral networks and locally connected networks on graphs,'' \emph{Proc. of the Intl. Conf. on Learning Representations}, Banff, AB, Canada, Apr. 2014.

%\bibitem{bronstein2017}
%M.M. Bronstein, J. Bruna, Y. LeCun, A. Szlam, and P. Vandergheynst, ``Geometric deep learning,'' \emph{IEEE Signal Processing Magazine}, vol. 34(4), pp. 18-42, Jul. 2017.

\bibitem{du2018}
J. Du, S. Zhang, G. Wu. J.M.F. Moura and S. Kar, ``Topology adaptive graph convolutional networks,'' \emph{arXiv:1710.10370v5}, Feb. 2018.

\bibitem{cheung2020}
M. Cheung, J. Shi, O. Wright, L.Y. Jiang, X. Liu, and J.M.F. Moura, ``Graph signal processing and deep learning: convolution, pooling and topology,'' vol. 37 (6), pp. 139-149. Nov. 2020.
     
\bibitem{gama2020}
F. Gama, E. Isufi, G. Leus, and A. Ribeiro, ``Graphs, convolutions, and neural networks,'' \emph{IEEE Signal Processing Magazine}, vol. 37 (6), pp. 128-138, Nov. 2020.

\bibitem{zhang2019}
S. Zhang, H. Tong, J. Xu, and R. Maciejewski, ``Graph convolutional networks: A comprehensive review,'' \emph{Computational Social Networks}, vol. 6(11), 2019.

\bibitem{wu2020}
Z. Wu, S. Pan, F. Chen, G. Long, C. Zhang, and P.S. Yu, ``A comprehensive survey on graph neural networks,'' \emph{IEEE Trans. on Neural Networks and Learning Systems}, vol. , pp. 1-21, Mar. 2020.

\bibitem{sahoo2018}
D. Sahoo, Q. Pham, J. Lu, S.C.H. Hoi, ``Online deep learning: Learning deep neural networks on the fly,'' \emph{Proc. of the 27th Intl. Joint Conf. on Artificial Intelligence}, 2018.

\bibitem{monga2021}
V. Monga, Y. Li, and Y.C. Eldar, ``Algorithm unrolling: Interpretable, efficient deep learning for signal and image processing,'' \emph{IEEE Signal Processing Magazine}, vol. 38(2), pp. 18-44, Feb. 2021.

\bibitem{soloman2019}
O. Soloman \emph{et al.}, ``Deep unfolded robust PCA with application to clutter suppression in ultrasound,'' \emph{IEEE Trans. on Medical Imaging}, vol. 39(4), pp. 1051-1063, Sep. 2019.

\bibitem{servetnyk2019}
M. Servetnyk and C.C. Fung, ``Distributed Joint Transmitter Design and Selection Using Augmented ADMM,'' \emph{Proc. of the IEEE Intl. Conf. on Speech, Acoustics and Signal Processing}, Brighton, UK, May 2019.

\bibitem{servetnyk2020}
M. Servetnyk, C.C. Fung, and Z. Han, ``Unsupervised federated learning for unbalanced data,'' \emph{Proc. of the IEEE Global Communications Conference}, Taipei, Taiwan, Dec. 2020.

% SCIP-SDP
\bibitem{scip_sdp}
SCIP-SDP, ``A mixed-integer semidefinite programming plugin for SCIP,'', \href{http://www.opt.tu-darmstadt.de/scipsdp/}{\textit{http://www.opt.tu-darmstadt.de/scipsdp/}}.


% Wireless InSite
\bibitem{insite}
Wireless InSite, 3D Wireless Prediction Software, \emph{\textcolor{blue}{https://www.remcom.com/wireless-insite-em-propagation-software}}.

% IEEE DataPort
\bibitem{dataport}
IEEE DataPort, \emph{\textcolor{blue}{http://ieee-dataport.org}}.

\bibitem{shuman2017}
D.I. Shuman, P. Vandergheynst, D. Kressner, and P. Frossard, ``Distributed signal processing via Chebyshev polynomial approximation,'' \emph{arXiv1111.5239v3}, Jul. 2017.
\end{comment}

% Multitarget Tracking 
%\bibitem{vo2015multitarget}
%B.-N. Vo \emph{et al.}, ``Multitarget tracking,'' \emph{Wiley Encyclopedia of Electrical and Electronics Engineering}, Sep. 2015.

%\bibitem{qiu2015survey}
%C. Qiu \emph{et al.}, ``A survey of motion-based multitarget tracking methods,'' \emph{Progress In Electromagnetics Research B}, vol. 62, pp. 195-223, 2015.

\end{thebibliography}

\end{document}




